\section{Forms, birational maps, and uniform inputs}
\label{sec:preliminaries}

The analytic and Hodge-theoretic arguments are over $\C$. We return to an
arbitrary algebraically closed characteristic-zero field at the end of
the proof. Varieties in the intervening sections are normal, integral,
and projective unless otherwise specified. We use the standard discrepancy
conventions for klt, dlt, and lc pairs \citep{Kollar2013}.
For normal varieties, $m(K_X+\Delta)\sim0$ means that this is an
\emph{integral principal divisor}. In particular it is Cartier. This
convention is stronger than rational linear triviality of the same
multiple.

\subsection{Trivializations as rational forms}

Let $F=\C(X)$ for an $n$-fold $X$. A rational $m$-canonical form is a
nonzero element of
$(\bigwedge^n\Omega_{F/\C})^{\otimes m}$. If $\eta$ is a rational
top form and $\theta=u\eta^{\otimes m}$, put
\[
                     \Div(\theta)=\Div(u)+m\Div(\eta).
\]
The divisor on the right is independent of the expression for $\theta$.
The equality $m(K_X+\Delta)\sim0$ is equivalent to the existence of
such a form with
\begin{equation}\label{eq:form-trivialization}
                         \Div(\theta)=-m\Delta.
\end{equation}
All forms satisfying this identity differ by a nonzero constant.
The identity also implies that $m\Delta$ is integral.

A birational map between pairs is \emph{$B$-birational}, or
\emph{crepant birational}, if the log canonical pullbacks agree on a
common resolution, with canonical divisors chosen compatibly. We write
$\Bir(X,\Delta)$ for the group of these self-maps. A
\emph{pseudo-automorphism} is a birational self-map which is an
isomorphism on open subsets with complements of codimension at least two.

\begin{lemma}\label{lem:form-crepant}
Let $(X,\Delta)$ and $(Y,\Delta_Y)$ be normal projective pairs, and let
$\theta_X,\theta_Y$ be degree-$m$ rational forms satisfying
\eqref{eq:form-trivialization} for the respective boundaries.
A birational map $f:X\dashrightarrow Y$ is $B$-birational if and only if
\[
                       f^*\theta_Y=c\theta_X
                 \quad\text{for some }c\in\C^*.
\]
\end{lemma}

\begin{proof}
Choose a common smooth projective resolution $W$, and use one rational
top form to represent $K_W$. Write the two pulled-back forms as
$u_X\eta^{\otimes m}$ and $u_Y\eta^{\otimes m}$. If
$K_W+\Delta_W^X$ and $K_W+\Delta_W^Y$ are the corresponding crepant
pullbacks, then
\[
             m(K_W+\Delta_W^X)=-\Div(u_X),\qquad
             m(K_W+\Delta_W^Y)=-\Div(u_Y).
\]
Their equality is equivalent to $\Div(u_Y/u_X)=0$. A rational function
with zero divisor on the normal projective integral variety $W$ is a
nonzero constant.
\end{proof}

Thus a $B$-birational self-map acts on the one-dimensional space of
trivializing forms by a scalar. Theorem~\ref{thm:characters} will bound
its order uniformly. No finiteness assumption on the order of the
birational map itself is made.

\subsection{The normal index and good-model inputs}

We isolate two substantial companion results. This makes clear which
uniformity is already available on normal varieties and which must be
proved for descent across the conductor.

\begin{theorem}[Normal index theorem {\citep[Theorem~1.2]{UPI}}]
\label{input:normal-index}
For every integer $d\ge0$ and every rational set
$I\subset[0,1]$ satisfying the descending chain condition, there is
an integer $a_{\mathrm{lc}}(d,I)>0$ such that every normal projective
integral lc pair $(X,\Delta)$ of dimension $d$, with coefficients in
$I$ and $K_X+\Delta\sim_{\Q}0$, satisfies
\[
                       a_{\mathrm{lc}}(d,I)(K_X+\Delta)\sim0.
\]
\end{theorem}

\begin{theorem}[Good minimal models {\citep[Theorem~11.1]{LA}}]
\label{input:good-model}
A projective lc pair over $\C$ with effective rational boundary and
pseudo-effective log canonical divisor has a good log minimal model.
In particular its adjoint is nonvanishing and its transform on the
good model is semiample.
\end{theorem}

For a klt pair with a good minimal model, an MMP with scaling may be
chosen to terminate at such a model. We use this consequence together
with the extraction, small $\Q$-factorialization, and big-adjoint
finite-generation results of \citet{BCHM2010}. When the adjoint is merely
pseudo-effective, the good-model theorem above is the nonvanishing
and semiampleness input.

We also use the uniform pluricanonical degree theorem
\citep[Theorem~1.1]{UPI}: for smooth projective varieties of each fixed
dimension and nonnegative Kodaira dimension, one complete pluricanonical
system defines the Iitaka fibration. Section~\ref{sec:rank} states the
chain and tracking consequences in the precise forms needed there.

\subsection{Bounded comparison models for rationally connected bases}

The rank argument uses a boundedness consequence of the canonical bundle
formula. Recall that a \emph{contraction} $h:V\to Y$ is a projective
surjective morphism of normal varieties with $h_*\OO_V=\OO_Y$.
Two normal varieties are isomorphic in codimension one if they have
isomorphic open subsets whose complements have codimension at least two.

\begin{proposition}\label{input:bounded-base}
Fix $n$. Let $h:V\to Y$ be a contraction over $\C$, where $V$ is
klt of dimension $n$, $K_V\sim_{\Q}0$, and $0<b=\dim Y<n$.
Suppose that a smooth projective resolution of $Y$ is rationally
connected. Then $Y$ is isomorphic in codimension one to a normal
projective variety $\widehat Y$ in a bounded family depending only
on $n$. In particular $\widehat Y$ has a very ample Cartier divisor
$H$ for which
\begin{equation}\label{eq:bounded-base-degrees}
            H^b\le C_n,\qquad
            \abs{K_{\widehat Y}\cdot H^{b-1}}\le C_n.
\end{equation}
The canonical intersection is the degree of the canonical Weil
divisor on a general complete-intersection curve in the smooth locus.
\end{proposition}

\begin{proof}
The bounded moduli denominator proposition of \citet[Proposition~4.1]{UPI}
applies to $h$. Its lower-dimensional normal-index hypotheses are
supplied by Theorem~\ref{input:normal-index}. It gives a generalized
klt pair $(Y,B_Y+M_Y)$ with effective boundary in a dimension-dependent
rational DCC set, b-nef moduli data with a uniformly bounded
b-Cartier denominator, and
\[
                         K_Y+B_Y+M_Y\sim_{\Q}0.
\]
Here the last equivalence follows from the exact canonical bundle
formula and $h_*\OO_V=\OO_Y$.
The proof of the rationally connected generalized torsion result
\citep[Proposition~4.5]{UPI} shows that these bases are uniformly
generalized $\epsilon$-lc and applies \citet[Theorem~1.7]{Birkar2025}
to obtain boundedness up to isomorphism in codimension one. Its normal
bounded comparison models give $\widehat Y$.

Choose very ample divisors from bounded projective embeddings of this
family. Their top degrees are bounded. A general curve cut out by
$b-1$ such divisors is smooth and avoids the singular locus, and its
degree and genus are bounded. Adjunction gives
\[
            K_{\widehat Y}\cdot H^{b-1}
                  =2g(C)-2-(b-1)H^b,
\]
which proves the second bound. For $b=1$ use the normal curve itself.
\end{proof}

\subsection{Generic data and section orders}

All finite collections of varieties, divisors, maps, and forms can be
defined over a finitely generated subfield of $\C$. Geometric generic
fibres and moving subvarieties may be studied after algebraically closed
field extension; sections commute with that extension, and the
singularity conditions can be checked on a descended log resolution.
We use very general points when countably many section-order conditions
must hold simultaneously. Equivalently, the points may be taken geometric
generic over a countable field containing the data. Dominant extensions
of parameter spaces used to define models preserve this genericity.

For a big rational divisor $D$ on an $n$-fold, write
$\vol(D)=\limsup_{m\to\infty}n!h^0(mD)/m^n$, with degrees divisible
enough to make the divisor integral. Divisors compared by an effective
rationally linearly equivalent difference give inclusions of these
complete section spaces. Section~\ref{sec:rank} will use such inclusions
on resolutions, where orders are measured at smooth generic points.
Object-dependent divisibility and asymptotic thresholds are permitted;
the asserted final bounds depend only on the parameters specified in
the corresponding statements.
